\documentclass[12pt,a4paper]{article}
\usepackage[utf8]{inputenc}
\usepackage{amsmath,amssymb}
\usepackage{graphicx}
\usepackage[left=1.5cm,right=1cm,top=1.3cm,bottom=1.3cm]{geometry}
\usepackage{lastpage}
% Font Times New Roman (cần XeLaTeX + fontspec — uncomment nếu VPS hỗ trợ):
% \usepackage{fontspec}
% \setmainfont{Times New Roman}

\begin{document}

\begin{center}
\begin{tabular}{p{0.4\textwidth}p{0.6\textwidth}}
\textbf{SỞ GDĐT SƠN LA} & ĐỀ THI CUỐI KÌ I NĂM 2026-2027 \tabularnewline
\textbf{TRƯỜNG THPT CÂY MƠ} & \textbf{Môn: Toán 12} \tabularnewline
\textbf{Đề chính thức} & \textit{Thời gian: 90 phút, không kể thời gian phát đề} \tabularnewline
\textit{(Đề thi gồm có .... trang)} & \tabularnewline
Họ và tên: ...................................................................................... \newline Số báo danh: .................................................................................. & \textbf{Mã đề: 9421} \tabularnewline
\end{tabular}
\end{center}
\vspace{0.5cm}

\noindent\textbf{PHẦN I. Câu trắc nghiệm nhiều phương án lựa chọn.}
Thí sinh trả lời từ câu 1 đến câu 2. Mỗi câu thí sinh chỉ chọn một phương án.

\medskip

\noindent\textcolor{blue}{\textbf{Câu 1.}} Cho hàm số $y = f(x)$ xác định trên $\mathbb{R}$ và có bảng biến thiên như hình vẽ sau

\begin{center}
\includegraphics[width=0.6\textwidth]{images/tikz_178c2b460ac4.svg}
\end{center}

Hàm số đã cho đồng biến trên khoảng nào dưới đây?

\noindent\begin{tabular}{p{0.25\linewidth} p{0.25\linewidth} p{0.25\linewidth} p{0.25\linewidth}}
\textbf{A.}~$(0;1)$. & \textbf{B.}~$(-\infty;-1)$. & \textbf{C.}~$(-1;1)$. & \textbf{D.}~$(-1;0)$. \\
\end{tabular}

\noindent\textcolor{blue}{\textbf{Câu 2.}} Cho hàm số $y=f(x)$ có bảng biến thiên như sau
	
\begin{center}
\includegraphics[width=0.6\textwidth]{images/tikz_0a1d3651e591.svg}
\end{center}

	Hàm số đã cho nghịch biến trên khoảng nào dưới đây?

\noindent\begin{tabular}{p{0.25\linewidth} p{0.25\linewidth} p{0.25\linewidth} p{0.25\linewidth}}
\textbf{A.}~$(1 ;+\infty)$. & \textbf{B.}~$(0 ; 1)$. & \textbf{C.}~$(-1 ; 0)$. & \textbf{D.}~$(0 ;+\infty)$. \\
\end{tabular}

\noindent\textbf{PHẦN II. Câu trắc nghiệm đúng sai.}
Thí sinh trả lời từ câu 1 đến câu 2. Trong mỗi ý, thí sinh chọn đúng (Đ) hoặc sai (S).

\medskip

\noindent\textcolor{blue}{\textbf{Câu 1.}} Cho hàm số $y=f(x)=\dfrac{ax^2+bx+c}{mx+n}$ (với $a, m \ne 0$) có đồ thị là đường cong như hình.

\begin{center}
\includegraphics[width=0.6\textwidth]{images/tikz_edc318573428.svg}
\end{center}


\noindent a)~Hàm số đồng biến trên các khoảng $(-\infty;-3)$ và $\left(-1;+\infty\right)$.

\noindent b)~$f(2024)<f(2025)$.

\noindent c)~Đồ thị hàm số có tiệm cận đứng $x=-3$ và đường tiệm cận xiên $y=x+1$.

\noindent d)~Đồ thị hàm số đi qua điểm $A\left(2;\dfrac{13}{4}\right)$.

\noindent\textcolor{blue}{\textbf{Câu 2.}} Trong thời gian $8$ phút đầu kể từ khi xuất phát, độ cao $h$ (tính bằng mét) so với mặt đất của một quả bóng bay vào thời điểm $t$ phút được cho bởi công thức $h(t) = 6t^3 - 81t^2 +324t$.

\noindent a)~Tại thời điểm ban đầu, quả bóng đang ở mặt đất.

\noindent b)~Sau $1$ phút, quả bóng bay đạt độ cao $250$ mét so với mặt đất.

\noindent c)~Trong $3$ phút đầu tiên, quả bóng bay tăng dần độ cao.

\noindent d)~Từ phút thứ $3$ đến phút thứ $8$, quả bóng bay đang tăng dần độ cao.

\noindent\textbf{PHẦN III. Câu trả lời ngắn.}
Thí sinh trả lời từ câu 1 đến câu 2.

\medskip

\noindent\textcolor{blue}{\textbf{Câu 1.}} Biết rằng khoảng nghịch biến lớn nhất của hàm số $y = \dfrac{x^2 + 2x + 2}{x + 1}$ là hai khoảng $(a;b)$, $(b;c)$ với $a < b < c$. Tính $T = a + b + c$.
\par

\noindent\textcolor{blue}{\textbf{Câu 2.}} Thể tích $V$ của $1$ kg nước (tính bằng cm$^3$) ở nhiệt độ $T$ (đơn vị $^\circ$C) khi thay đổi từ $0^\circ$C đến $30^\circ$C được cho xấp xỉ bởi công thức
	$V=999,87-0,06426T+0,0085043T^2-0,0000769T^3$ (nguồn: James Stewart, T. ($2015$). Calculus.Cengage Learning $8$th edition, p.$284$). Biết rằng kể từ $T_0^\circ$ C trở đi, với $0^\circ\le T_0\le 30^\circ$ thể tích $V$ tăng, tìm $T_0$ (kết quả làm tròn đến hàng đơn vị).

\noindent\textbf{PHẦN IV. Câu tự luận.}
Thí sinh trả lời từ câu 1 đến câu 2.

\medskip

\noindent\textcolor{blue}{\textbf{Câu 1.}} [Trích đề ôn tập HKI-Lớp 12-THPT Bùi Thị Xuân-Năm học 2024-2025]
	Một vật đang đứng yên thì bắt đầu chuyển động theo quy luật $s(t)=-t^3+3t^2+6t$, với $t$ (giây) là khoảng thời gian tính từ lúc vật bắt đầu chuyển động và $s$ (mét) là quãng đường vật đi được trong khoảng thời gian đó. Hỏi vật tăng tốc trong khoảng thời gian bao nhiêu giây tính từ lúc bắt đầu chuyển động?

\noindent\textcolor{blue}{\textbf{Câu 2.}} Sau khi phát hiện một bệnh dịch, các chuyên gia y tế ước tính số người nhiễm bệnh kể từ ngày xuất hiện bệnh nhân đầu tiên đến ngày thứ $t$ là $f(t)=45t^2-t^3$, $t=0$, $1$, $2$, $\ldots$, $25$. Nếu coi $f(t)$ là hàm số xác định trên đoạn $[0; 25]$ thì đạo hàm $f'(t)$ được xem là tốc độ truyền bệnh (người/ngày) tại thời điểm $t$. Giả sử khoảng thời gian mà tốc độ truyền bệnh giảm là từ ngày thứ $m$ đến ngày thứ $n$. Khi đó $n-m$ bằng bao nhiêu?


\textbf{------------ HẾT ------------}
\label{lastpage}

\end{document}